% !TEX root = ../main.tex
\chapter{Treball previ i evolució del pipeline C}\label{ch:pipeline}
\fancyhead[RE]{\slshape \textsf{Capítol \thechapter.\ Pipeline C}}
\fancyhead[LO]{\slshape \textsf{\nouppercase \rightmark}}

\section{Antecedent propi}\label{sec:antecedent}

El punt de partida d'aquesta memòria és un treball propi anterior sobre la
portabilitat de SORTENY a C \cite{anez2026treball}. Aquell treball ja havia
demostrat que una part substancial del compressor neuronal podia executar-se
sense el runtime de TensorFlow: es van exportar pesos, es van implementar
operadors de la transformada d'anàlisi i es va reduir el consum de memòria amb
buffers reutilitzables.

En una imatge de prova, el compressor C anterior va passar d'uns 796,5~MiB de
la ruta Python a uns 87~MiB, i va obtenir una acceleració aproximada
d'$1,47\times$. La comparació dels latents va mostrar un 99,9876\% de valors
idèntics, amb discrepàncies de $\pm1$ en una fracció petita. Aquest resultat
era suficient per validar la direcció, però encara no constituïa un còdec
complet: faltava una ruta de descompressió, un mecanisme de qualitat espacial,
una validació de bitstream i una mesura de bitrate codificat.

\subsection{Què s'hereta i què s'afegeix}

La Taula~\ref{tab:herencia} separa el treball heretat de les aportacions
d'aquesta memòria. La lectura important és que no es torna a entrenar SORTENY:
es construeix una ruta C al voltant del model entrenat i se'n completa el cicle
compressió--descompressió.

\begin{table}[H]
\centering
\scriptsize
\begin{tabularx}{\textwidth}{p{2.9cm}X X X}
\toprule
Component & Estat al treball previ & Aportació actual & Per què importa \\
\midrule
Pesos entrenats & Exportació inicial de pesos de l'encoder. & Validació de
formes, pesos d'encoder i decoder i índexs de càrrega. & Evita errors de
layout que poden donar sortides plausibles però incorrectes. \\
Compressor C & Transformada d'anàlisi funcional. & Integració amb informació
de qualitat local i contenidor autocontingut. & Permet transportar la qualitat
aplicada juntament amb els latents. \\
Descompressor & No era la ruta principal. & Implementació completa de síntesi,
validació de capçalera i reconstrucció RAW. & Tanca el còdec i permet verificar
que el fitxer generat és llegible. \\
Qualitat & Un únic paràmetre global. & Mapa local de qualitat i calibració
associada. & Fa possible protegir unes zones i degradar-ne unes altres. \\
Bitrate & Proxies sobre latents. & Mesura externa amb codificació real de la
referència. & Separa compressibilitat estimada de bitrate codificat. \\
Raspberry & Prova inicial de compressor. & Pipeline complet, optimització i
mesura de cost de qualitat local. & Aproxima el comportament en hardware de
baix recurs. \\
\bottomrule
\end{tabularx}
\caption{Separació entre l'antecedent propi i les aportacions d'aquesta
memòria.}
\label{tab:herencia}
\end{table}

\section{Referència Python i contracte funcional}\label{sec:python_ref}

SORTENY disposa d'una implementació Python de referència. En aquesta memòria
no es presenta com a producte operatiu, sinó com a definició executable del
model entrenat: conté la càrrega de checkpoints, les transformades, la
modulació de qualitat i el codificador final de TensorFlow Compression. Els
camins exactes del repositori i les ordres de reproducció es documenten als
apèndixs i al repositori públic associat; en el cos principal només interessa
la seva funció metodològica.

Portar aquesta referència a C no consisteix a traduir sintaxi. En Python, el
framework conserva molta informació de forma implícita: ordre de dimensions,
padding de convolucions, disposició dels kernels, tipus de dades, arrodoniment,
normalització i empaquetat final. En C, totes aquestes decisions han de quedar
fixades de manera explícita. Un error petit en un índex pot no provocar cap
fallada d'execució, però sí canviar els latents i, per tant, la reconstrucció.

La referència Python manté dues responsabilitats durant aquest treball:

\begin{itemize}[leftmargin=*]
  \item servir com a patró funcional del model entrenat;
  \item proporcionar una codificació final real, mitjançant fitxers
  \code{.tfci}, per mesurar bitrate quan el contenidor C encara no inclou un
  codificador per entropia o aritmètic equivalent.
\end{itemize}

\subsection{Frontera entre referència i ruta C}

La Taula~\ref{tab:python_c_contract} resumeix la frontera. El criteri és
pragmàtic: la ruta C ha d'executar el model i reconstruir la imatge amb recursos
limitats; la referència conserva les eines de model i de codificació final que
encara no s'han portat a C.

\begin{table}[H]
\centering
\small
\begin{tabularx}{\textwidth}{p{3.3cm}X X}
\toprule
Funció & Referència de model & Ruta C operativa \\
\midrule
Càrrega del model & Checkpoints i objectes TensorFlow. & Fitxers de pesos
binaris i índexs de formes. \\
Qualitat global & Valor continu del paràmetre de qualitat. & Valor global
equivalent quantitzat o transmès com a configuració. \\
Qualitat local & Array local acceptat per la referència. & Mapa local de
1.024 bytes transportat dins el contenidor C. \\
Transformades & Operadors del framework. & Bucles C, OpenMP i buffers
explícits. \\
Codificació final & TensorFlow Compression, fitxer \code{.tfci}. & Latents
enters dins un contenidor intermedi, sense codificació final eficient. \\
Ús al projecte & Patró funcional i bitrate real. & Ruta candidata per a
execució en hardware limitat. \\
\bottomrule
\end{tabularx}
\caption{Contracte entre la referència de model i la ruta C.}
\label{tab:python_c_contract}
\end{table}

La separació també condiciona la validació. Les diferències C--Python poden
venir de coma flotant, arrodoniment o detalls interns del framework. En canvi,
dues versions C comparades sota el mateix entorn han de produir la mateixa
reconstrucció si només s'ha optimitzat el codi. Per això la paritat C--C és el
criteri estricte quan es canvia la implementació C, mentre que la referència
Python actua com a comprovació externa del comportament del model.

\section{Implementació C del model}\label{sec:c_software_arch}

El model entrenat es considera congelat. La ruta C no aprèn pesos nous ni
modifica l'arquitectura; implementa la inferència, la lectura de pesos, la
informació de qualitat local i el contenidor necessari perquè el decoder pugui
reconstruir. Aquesta separació és important: les decisions científiques del
model provenen de SORTENY, mentre que aquesta memòria se centra en fer-lo
executables, verificable i mesurable en un entorn limitat.

\subsection{Mòduls i responsabilitats}\label{sec:pesos}

La implementació es divideix en responsabilitats conceptuals. Els noms exactes
de fitxer són traçables al repositori, però la Taula~\ref{tab:c_modules} mostra
la idea d'arquitectura.

\begin{table}[H]
\centering
\small
\begin{tabularx}{\textwidth}{p{3.0cm}X X}
\toprule
Mòdul conceptual & Responsabilitat & Validació associada \\
\midrule
Càrrega de model & Llegeix pesos, formes i índexs. & Mides esperades,
coherència de canals i presència de pesos mínims. \\
Operadors neuronals & Executa convolucions, GDN/IGDN, capes denses i
reordenacions. & Tests amb tensors petits i comparació amb valors esperats. \\
Entrada i sortida RAW & Llegeix BSQ \code{uint16}, controla endianess i escriu
la reconstrucció. & Mida exacta del RAW i rang $[0,65535]$. \\
Compressor & Aplica transformades, modulació i quantització. & Contenidor
estructuralment vàlid i latents coherents. \\
Descompressor & Valida el contenidor, desfà la modulació i sintetitza la
imatge. & Reconstrucció completa i verificació de dimensions. \\
Generadors de qualitat & Calculen el mapa local de qualitat a partir de
configuració o criteris automàtics. & Mapa de 1.024 bytes i traçabilitat de la
decisió. \\
\bottomrule
\end{tabularx}
\caption{Responsabilitats principals de la ruta C.}
\label{tab:c_modules}
\end{table}

Els pesos no es carreguen com una col·lecció opaca de fitxers. Cada tensor va
acompanyat d'un índex que indica nom, tipus, forma i mida. Aquesta informació
és fonamental perquè la xarxa combina moltes dimensions: canals d'entrada,
canals de sortida, kernels espacials, factors de modulació i transformades
espectrals. Si un kernel es llegís transposat, el programa podria acabar sense
errors però generaria una imatge incorrecta. Per això el carregador comprova la
mida esperada i la compatibilitat de canals abans d'associar cada buffer al
model.

\subsection{Punts crítics de la traducció a C}\label{sec:operadors}

Els operadors C reprodueixen la mecànica descrita al
Capítol~\ref{ch:fonaments}: convolucions amb stride, GDN i IGDN, capes denses
del modulador, reordenacions espacials i transformada espectral. El problema no
és escriure una fórmula a C, sinó conservar totes les convencions que el
framework gestiona de forma implícita.

Hi ha quatre punts especialment sensibles. Primer, les convolucions redueixen
resolució i canvien el nombre de canals; una transposició de kernel pot produir
una sortida amb la mida correcta però amb valors incorrectes. Segon, GDN i IGDN
fan sumes entre canals amb paràmetres apresos $\beta$ i $\gamma$; canviar
l'ordre o confondre canals altera la normalització. Tercer, el modulador de
qualitat genera factors que escalen els latents abans de l'arrodoniment; per
tant, un error de qualitat es propaga directament al fitxer. Quart, les
reordenacions del decoder no fan operacions de coma flotant, però són molt
sensibles a l'ordre de canals i píxels.

Aquests detalls expliquen per què els tests d'operadors no són un accessori:
un error de layout pot quedar poc visible en una imatge reconstruïda, però
trencar la paritat i fer que el decoder llegeixi una estructura diferent de la
que el compressor ha escrit.

La transformada espectral també té dues cares: anàlisi abans de la compressió
espacial i síntesi després del decoder. La ruta oficial utilitza pesos de
síntesi validats. El codi conserva una salvaguarda de compatibilitat per
derivar una inversa quan falten aquests pesos, però aquesta ruta no és la base
de cap resultat experimental de la memòria.

\subsection{Per què desacoblar qualitat i inferència}

Una decisió de disseny és separar la generació de qualitat local de la
inferència del model. Un bloc de codi decideix quina qualitat ha de tenir cada
cel·la de la graella; un altre bloc executa les transformades neuronals amb
aquesta informació. La frontera entre tots dos blocs és el mapa de qualitat:
una graella de 1.024 bytes que el compressor només consumeix, sense saber si
prové d'una qualitat global, d'una calibració o d'un criteri automàtic.

Aquesta separació redueix risc experimental. Afegir un criteri nou de selecció
de regions no modifica les convolucions ni els pesos; només canvia el mapa que
entra al compressor. Inversament, optimitzar GDN, IGDN o la transformada
espectral no altera la decisió de quina regió és prioritària. Així es poden
provar polítiques de qualitat i optimitzacions C per separat, amb validacions
diferents.

Aquest desacoblament és el pont cap al Capítol~\ref{ch:control}: aquí només es
descriu que el compressor accepta una qualitat global o local; el capítol
següent defineix formalment quins modes de qualitat s'han implementat.

\section{Gestió de memòria}\label{sec:memoria}

La memòria és una restricció central. Una implementació directa podria reservar
la sortida de cada capa per a totes les bandes i conservar-la fins al final.
Això és còmode en un framework de xarxa neuronal, però no és adequat per a una
placa de baix recurs. La ruta C tracta els tensors intermedis com a dades amb
vida curta: quan una capa ja no necessita una sortida anterior, aquell espai es
pot reutilitzar.

La idea principal és doble:

\begin{itemize}[leftmargin=*]
  \item utilitzar buffers contigus i reutilitzables per a les activacions més
  grans;
  \item processar el flux per bandes latents, en lloc de mantenir totes les
  bandes i tots els tensors intermedis simultàniament.
\end{itemize}

Per una banda, el compressor normalitza el pla d'entrada, executa els nivells
d'anàlisi, aplica la modulació de qualitat, arrodoneix els latents i els escriu
al contenidor abans de passar a la banda següent. El decoder segueix la idea
inversa: llegeix una banda latent, desfà la modulació, executa la síntesi i
només conserva el tensor espectral final necessari per reconstruir el RAW.

La diferència és qualitativa. Sense streaming, els latents intermedis de les
vuit bandes i les sortides de capes grans competirien per memòria al mateix
temps. Amb streaming, el pic queda dominat pels buffers més grans i pel conjunt
de pesos, no per una còpia completa de cada etapa per cada banda.

\subsection{Pressupost nominal de buffers}

Les mides de la Taula~\ref{tab:memory_budget} són nominals: es calculen a
partir de la forma del tensor i del tipus de dada, no són mesures de RSS. Per
exemple, un tensor \code{float32} amb $128\cdot256\cdot256$ elements ocupa
$128\cdot256\cdot256\cdot4$ bytes, és a dir, uns 32~MiB. Les mesures reals de
Raspberry es presenten al Capítol~\ref{ch:bord}.

\begin{table}[H]
\centering
\small
\begin{tabular}{lrrr}
\toprule
Reserva principal & Forma típica & Tipus & Mida nominal \\
\midrule
Entrada planificada & $8\times512\times512$ & \code{float32} & 8,0~MiB \\
Tensor espectral & $8\times512\times512$ & \code{float32} & 8,0~MiB \\
Buffer intermedi A & $128\times256\times256$ & \code{float32} & 32,0~MiB \\
Buffer intermedi B & $128\times256\times256$ & \code{float32} & 32,0~MiB \\
Latent d'una banda & $384\times32\times32$ & \code{float32} & 1,5~MiB \\
Cache de moduladors & $256\times3072$ & \code{float32} & 3,0~MiB \\
Pesos encoder & formes de l'índex & \code{float32} & 10,3~MiB \\
\bottomrule
\end{tabular}
\caption{Ordre de magnitud de les reserves principals. Les vides útils se
solapen parcialment i el RSS no és la suma aritmètica de totes les files.}
\label{tab:memory_budget}
\end{table}

La Taula~\ref{tab:memory_budget} mostra per què la reutilització és necessària:
només dos buffers intermedis ja representen desenes de MiB. Duplicar aquests
buffers per paral·lelitzar bandes completes augmentaria ràpidament la pressió
de memòria. Per això el paral·lelisme s'aplica dins d'operadors, píxels o
canals, mentre el flux general manté controlat el nombre de tensors grans vius
alhora.

La cache de moduladors és una excepció deliberada. Ocupa uns pocs MiB, però
evita recalcular la xarxa auxiliar de qualitat per cada cel·la quan molts blocs
comparteixen el mateix valor. És una decisió de compromís: gastar una petita
quantitat de memòria per fer més predictible el temps de compressió.

\section{Cicle compressor--contenidor--descompressor}
\label{sec:compressor_c}

El pipeline C complet es pot llegir com un cicle únic:

\begin{figure}[H]
\centering
\fbox{\begin{minipage}{0.92\textwidth}
\centering
RAW BSQ \code{uint16}
$\rightarrow$ anàlisi espectral
$\rightarrow$ anàlisi espacial
$\rightarrow$ modulació de qualitat
$\rightarrow$ latents enters
$\rightarrow$ contenidor C
$\rightarrow$ descompressió
$\rightarrow$ RAW reconstruït
\end{minipage}}
\caption{Flux funcional del còdec C implementat.}
\label{fig:pipeline_encoder}
\end{figure}

El compressor rep la imatge, els pesos i la informació de qualitat. Si la
qualitat és global, tots els blocs comparteixen el mateix valor. Si és local,
cada cel·la de la graella utilitza el byte corresponent. La definició exacta
d'aquests modes queda reservada al Capítol~\ref{ch:control}; aquí només importa
que el valor de qualitat viatja juntament amb els latents perquè el decoder
pugui desfer la modulació aplicada.

El descompressor és una aportació essencial d'aquesta memòria. Llegeix el
contenidor, valida la capçalera, comprova dimensions, carrega els pesos de
síntesi, reconstrueix els moduladors necessaris, desfà la modulació dels
latents, aplica les capes de síntesi, recupera les bandes i escriu el RAW
final. Si la sortida no té la mida esperada, la reconstrucció es considera
invàlida encara que el procés hagi acabat.

\subsection{Format del contenidor C}\label{sec:bitstream_c}

El format actual és deliberadament simple. No és un contenidor CCSDS ni un
format de vol. La disposició és:

\begin{table}[H]
\centering
\small
\begin{tabular}{lrrl}
\toprule
Camp & Tipus & Elements & Descripció \\
\midrule
Capçalera & \code{uint16} & 5 & Bandes, alçada, amplada, datatype i filtres. \\
Mapa de qualitat & \code{uint8} & 1.024 & Qualitat per cada cel·la latent. \\
Latents & \code{int32} & $8\cdot384\cdot32\cdot32$ & Ordre banda, canal,
fila, columna. \\
\bottomrule
\end{tabular}
\caption{Disposició del contenidor intermedi C.}
\label{tab:bitstream}
\end{table}

La capçalera ocupa 10 bytes i el mapa de qualitat ocupa sempre 1.024 bytes. En
la implementació actual, fins i tot una qualitat global fixa s'escriu com un
mapa constant. Això és ineficient en sentit estricte, perquè es podria guardar
un únic valor i un flag, però simplifica el contracte del decoder: sempre llegeix
la mateixa estructura i sempre sap quina qualitat correspon a cada cel·la.

El cost d'aquesta decisió és petit davant dels latents enters. 1.024 bytes són
8.192 bits, és a dir:

\begin{equation}
\frac{8.192}{8\cdot512\cdot512}=0,00390625\;\mathrm{bps}.
\end{equation}

Els latents, en canvi, ocupen una mida gairebé fixa perquè encara no s'aplica
una codificació final eficient. Per tant, la mida d'aquest contenidor no pot
demostrar l'estalvi de bits d'una política de qualitat. La seva funció és una
altra: transportar tot el que el decoder necessita per reconstruir i permetre
auditar la ruta C.

Aquest contenidor tampoc pretén ser el format final d'una missió. Un format
operatiu hauria d'afegir metadades de versió, comprovació d'integritat i una
codificació eficient dels latents. Aquí es conserva una estructura simple per
centrar la validació en el còdec i en la qualitat aplicada.

\subsection{Mida i implicacions}

Per una imatge $8\times512\times512$, el nombre d'enters latents és:

\begin{equation}
N_y=8\cdot384\cdot32\cdot32=3.145.728.
\end{equation}

Els quatre factors provenen directament de la geometria interna: vuit bandes,
384 canals latents per banda i una graella espacial latent de $32\times32$.
Cada valor latent es guarda com un enter de 32 bits, és a dir, quatre bytes.
Emmagatzemats com \code{int32}, ocupen 12.582.912 bytes. Afegint capçalera i
mapa de qualitat, el contenidor C és més gran que el RAW de 4.194.304 bytes.
No és una contradicció: és una representació intermèdia per validar
transformada, quantització, transport de qualitat i descompressió. El guany de
compressió només apareix quan la distribució dels latents s'explota amb un
model probabilístic i un codificador final.

Per aquest motiu, els proxies com entropia o gzip només serveixen per orientar
la selecció de candidats. Poden indicar que uns latents semblen més fàcils de
codificar, però no substitueixen la mida \code{.tfci} obtinguda amb la
referència. La metodologia del Capítol~\ref{ch:metodologia} explica com
s'utilitza aquesta referència externa per mesurar bitrate real.

\subsection{Validació d'entrades}\label{sec:decoder_c}

El descompressor rebutja capçaleres incompletes, dimensions zero, nombre de
bandes no suportat, mides no divisibles per la graella latent, mapes de
qualitat truncats, latents incomplets, bytes sobrants i pesos incoherents. Són
comprovacions defensives: una dimensió corrupta podria provocar reserves
desproporcionades o accessos fora de límits. El format encara no incorpora una
comprovació d'integritat forta, de manera que detecta inconsistència
estructural, però no tota corrupció silenciosa que mantingui les mateixes mides.

\section{Validació del pipeline C}\label{sec:tests_c}

La validació no es basa en una única prova. Un PSNR global correcte pot amagar
un error petit de disposició; un test d'operador no garanteix que el fitxer
complet sigui llegible; una execució local no prova que el consum de memòria
sigui assumible a Raspberry. Per això s'utilitzen nivells complementaris.

\begin{table}[H]
\centering
\small
\begin{tabularx}{\textwidth}{p{3.2cm}X X}
\toprule
Nivell & Què comprova & Què no garanteix \\
\midrule
Operadors aïllats & Fórmules, índexs i layout en tensors petits. & Que el
pipeline complet gestioni bé fitxers i pesos. \\
Integració C & Compressió seguida de descompressió amb un contenidor llegible.
& Que el resultat coincideixi amb la referència de model. \\
Paritat entre versions C & Que una optimització no canvia Q, latents o RAW. &
Que el model C sigui idèntic a Python en cada operació interna. \\
Comparació amb referència & Que la semàntica general del model es conserva. &
Bit-exactitud absoluta en coma flotant. \\
Validació externa de bitrate & Que la qualitat sol·licitada arriba al
codificador real de la referència. & Que el contenidor C ja sigui un bitstream
final eficient. \\
Raspberry & Dependències, RSS, temperatura i temps en hardware limitat. & Una
certificació de vol o una garantia de temps real. \\
\bottomrule
\end{tabularx}
\caption{Nivells de validació del pipeline C.}
\label{tab:test_matrix}
\end{table}

La paritat entre versions C és especialment estricta. Si es canvia una
implementació per accelerar-la, la reconstrucció ha de romandre igual. Això
permet atribuir qualsevol millora de temps al canvi de codi i no a una
relaxació de qualitat. La referència Python continua sent necessària, però es
fa servir amb una altra finalitat: comprovar que la ruta C segueix representant
el model entrenat i obtenir, quan cal, una mesura de bitrate codificat.

\section{Límit del pipeline C actual}\label{sec:limit_c}

El pipeline C actual resol inferència, qualitat local, contenidor i
reconstrucció. Encara no resol la codificació final eficient. Els latents
s'emmagatzemen com enters dins un contenidor intermedi, sense model
probabilístic ni codificador per entropia o aritmètic equivalent al de la
referència. Això té dues conseqüències:

\begin{enumerate}[leftmargin=*]
  \item la mida del contenidor C és pràcticament independent del guany que
  podria obtenir un codificador final;
  \item el bitrate real s'ha de mesurar amb un codificador extern de referència
  fins que aquesta etapa es porti a C.
\end{enumerate}

Aquesta separació evita una conclusió incorrecta. El contenidor C demostra que
la ruta operativa pot transportar i reconstruir la qualitat aplicada; els
fitxers codificats de la referència demostren quin bitrate s'obté quan els
latents passen per una codificació final real. El treball futur natural és
integrar també aquesta última etapa a C.

La calibració exacta de l'escala de qualitat i els modes que generen el mapa es
defineixen al Capítol~\ref{ch:control}. Aquest capítol només estableix el
contracte que els fa possibles: la qualitat aplicada ha de viatjar amb els
latents i ha de poder ser llegida pel decoder.
